\documentclass[a4paper]{article}

\usepackage[fontset=fandol]{ctex}
\usepackage{amsmath, amssymb}
\usepackage{graphicx}
\usepackage[margin=2.45cm]{geometry}
\usepackage[most]{tcolorbox}
\usepackage{hyperref}
\usepackage{booktabs}
\usepackage{float}
\usepackage{array}
\usepackage{xcolor}
\usepackage{enumitem}
\usepackage{caption}

\setlength{\parindent}{2em}
\setlength{\parskip}{0.35em}
\setlength{\emergencystretch}{3em}
\linespread{1.06}
\sloppy
\raggedbottom
\vfuzz=4pt

\newcolumntype{L}[1]{>{\raggedright\arraybackslash}p{#1}}
\newcolumntype{C}[1]{>{\centering\arraybackslash}p{#1}}

\newcommand{\E}{\mathbb{E}}
\newcommand{\Dtrain}{\mathcal{D}^{\mathrm{train}}}
\newcommand{\Dtest}{\mathcal{D}^{\mathrm{test}}}
\newcommand{\Task}{\mathcal{T}}
\newcommand{\flearner}{F_{\phi}}
\newcommand{\classifier}{f_{\theta}}
\newcommand{\metaloss}{L(\phi)}
\newcommand{\taskloss}{\ell}

\hypersetup{
    colorlinks=true,
    linkcolor=blue!60!black,
    urlcolor=blue!60!black
}

\setlist[itemize]{leftmargin=2.2em,itemsep=0.22em,topsep=0.25em}
\setlist[enumerate]{leftmargin=2.4em,itemsep=0.22em,topsep=0.25em}
\setlist[description]{leftmargin=3.1em,itemsep=0.18em,topsep=0.25em}

\newtcolorbox{knowledgebox}[1]{
    enhanced, colback=blue!5!white, colframe=blue!75!black, colbacktitle=blue!75!black,
    coltitle=white, fonttitle=\bfseries, title={#1},
    attach boxed title to top left={yshift=-2mm, xshift=2mm},
    boxrule=1pt, sharp corners
}

\newtcolorbox{importantbox}[1]{
    enhanced, colback=yellow!10!white, colframe=yellow!80!black, colbacktitle=yellow!80!black,
    coltitle=black, fonttitle=\bfseries, title={#1}, sharp corners
}

\newtcolorbox{warningbox}[1]{
    enhanced, colback=red!5!white, colframe=red!75!black, colbacktitle=red!75!black,
    coltitle=white, fonttitle=\bfseries, title={#1}, sharp corners
}

\newcommand{\notetitle}{元学习（一）：Learning to Learn}
\newcommand{\notesubtitle}{把机器学习的三步骤提升到“学习算法本身”的层次}
\newcommand{\noteauthors}{基于李宏毅老师《机器学习 2025》课程内容整理}
\newcommand{\notedate}{\today}
\newcommand{\videochannel}{李宏毅深度学习课堂（Bilibili）}
\newcommand{\videoduration}{约 46 分 20 秒}
\newcommand{\videourl}{https://www.bilibili.com/video/BV1TAtwzTE1S?p=55}

\newcommand{\framefig}[4]{%
\begin{figure}[H]
    \centering
    \includegraphics[width=#1\textwidth]{#2}
    \caption{#3\protect\footnotemark}
\end{figure}
\footnotetext{视频画面时间区间：#4。}
}

\begin{document}
\pagestyle{empty}

\begin{center}
    \vspace*{1.0cm}
    {\Huge\bfseries \notetitle\par}
    \vspace{0.35cm}
    {\LARGE \notesubtitle\par}
    \vspace{0.75cm}
    \includegraphics[width=0.62\textwidth]{video.jpg}\par
    \vspace{0.75cm}
    {\large \noteauthors\par}
    {\large \notedate\par}
    \vspace{0.75cm}
    \begin{tcolorbox}[width=0.86\textwidth,center,sharp corners,
                      colback=black!3!white,colframe=black!50!white]
        \centering
        \textbf{视频信息}\\[3pt]
        频道：\videochannel \\
        时长：\videoduration \\
        链接：\href{\videourl}{\videourl}
    \end{tcolorbox}
\end{center}

\newpage
\tableofcontents
\newpage
\pagestyle{plain}

% =====================================================================
\section{为什么需要元学习}
% =====================================================================

Meta Learning 常被翻译为“元学习”，也常被概括成 learning to learn。课程一开始先解释这个主题为什么自然：很多机器学习实践并不只是“把数据喂给模型然后结束”，而是充满大量人工选择，例如网络架构、学习率、优化器、正则化强度、数据增强方式、训练轮数等。普通深度学习系统看起来很自动，实际上常常需要人反复试超参数。

\framefig{0.84}{figures/fig_01_hyperparameter_search.jpg}
{课程用“调参作业”引入元学习：深度学习系统的表现高度依赖网络结构、学习率等 hyperparameters；如果这些选择长期依赖人工经验，就说明“学习如何学习”本身也可以被系统化。}
{00:02:00--00:02:15}

传统做法中，研究者或工程师会根据经验决定一个学习算法的大量细节，再让这个算法从训练数据中学出模型。元学习反过来问：能否让机器也参与决定“如何学”？也就是说，不只让模型学参数，还让系统学习产生模型的程序、策略或初始化。

\begin{knowledgebox}{“Meta” 在这里不是玄学}
    这里的 meta 不是指模型比普通模型更神秘，而是指研究对象提高了一个层级。普通机器学习找的是某个任务上的函数 $f$，元学习找的是一个能够产生函数的学习算法 $F$。如果普通模型回答“这张图是不是猫”，元学习关心的是“给我少量训练例子以后，怎样快速得到一个能回答问题的分类器”。
\end{knowledgebox}

从工程角度看，元学习的动机至少有三类。
\begin{itemize}
    \item \textbf{减少手工设计：}把网络结构、初始化、学习率、更新规则等原本靠人工或网格搜索决定的部分，改成可以从任务经验中学习。
    \item \textbf{提升少样本适应：}希望系统看到少量新任务样本后，就能快速得到可用模型。
    \item \textbf{复用跨任务经验：}如果过去见过许多任务，就不应该每遇到一个新任务都从零开始摸索学习方式。
\end{itemize}

\begin{warningbox}{不要把元学习等同于某一个具体算法}
    元学习不是 MAML、Matching Networks 或 few-shot classification 的同义词。它是一种问题设定：学习对象从“任务函数”上升为“学习算法”。不同论文可以学习初始化、学习率、优化器、距离函数、任务表示、网络结构，甚至学习整个更新规则。
\end{warningbox}

\subsection{本章小结}

P55 的起点是调参困境：深度学习虽然自动学习权重，但大量“怎样学习”的选择仍然依赖人。元学习把这些选择放进可学习对象中，目标是让系统从一批任务中归纳出更好的学习算法。

% =====================================================================
\section{从普通机器学习三步骤到元学习三步骤}
% =====================================================================

课程用 Machine Learning 101 复习普通监督学习的三步骤：
\begin{enumerate}
    \item 定义一个带未知参数的函数集合，例如神经网络 $f_{\theta}$。
    \item 定义损失函数 $L(\theta)$，衡量某组参数在训练资料上的好坏。
    \item 用优化方法找出让损失小的参数 $\theta^\ast$。
\end{enumerate}

\framefig{0.84}{figures/fig_02_ml_three_steps.jpg}
{普通机器学习三步骤：先定义带未知参数的函数集合，再定义 loss function，最后用优化方法找出最佳参数；图中以猫狗分类网络为例，$\theta$ 表示神经网络中的可学习权重和偏置。}
{00:05:45--00:06:00}

这套三步骤的关键是：未知对象是函数 $f_{\theta}$ 的参数 $\theta$。训练完成后得到的是一个具体函数 $f_{\theta^\ast}$，它可以对测试输入做预测。元学习把这个结构向上提升一层：未知对象不再只是分类器参数，而是学习算法本身。

\framefig{0.84}{figures/fig_03_meta_learning_function.jpg}
{元学习也沿用机器学习的三步骤，只是要找的函数变成学习算法 $F$：它接收某个任务的训练样本，输出一个可用于该任务的分类器 $f^\ast$。}
{00:09:15--00:09:30}

用符号写，可以把普通机器学习看成
\[
    \Dtrain \xrightarrow{\text{hand-crafted learning algorithm}} f_{\theta^\ast}.
\]
其中学习算法通常由人设计，例如“用某个网络架构、交叉熵、Adam、固定学习率训练若干 epoch”。元学习则试图学习一个带参数的学习算法：
\[
    \flearner(\Dtrain) = f^\ast.
\]

符号说明如下。
\begin{description}
    \item[$\Dtrain$] 某个任务提供的少量训练样本。
    \item[$\flearner$] 带可学习组件 $\phi$ 的学习算法。
    \item[$f^\ast$] 由学习算法针对当前任务产生的模型，例如分类器。
    \item[$\phi$] 学习算法内部待学习的部分，可以是初始化、度量、更新规则、网络结构等。
\end{description}

\begin{importantbox}{层级差异}
    普通机器学习学习的是模型参数 $\theta$；元学习学习的是学习算法参数 $\phi$。$\theta$ 是某一个具体任务的模型参数，$\phi$ 则决定“看到任务训练样本以后如何产生 $\theta$ 或如何产生分类器”。这两个层级混在一起，后面的 support/query set、inner/outer loop 就会很容易看乱。
\end{importantbox}

\subsection{本章小结}

普通机器学习和元学习共享同一个形式：定义候选空间、定义评价函数、做优化。区别在于候选对象不同。普通机器学习的候选对象是任务函数 $f_{\theta}$，元学习的候选对象是学习算法 $F_{\phi}$。

% =====================================================================
\section{Step 1：学习算法里什么可以被学}
% =====================================================================

元学习的第一步是定义“学习算法的集合”。课程用问题 What is learnable in a learning algorithm? 来展开：一个学习算法并不只是抽象动词“训练”，它包含许多组件。我们可以规定哪些组件由人手工设定，哪些组件交给数据来学。

\framefig{0.84}{figures/fig_04_learnable_components_phi.jpg}
{Step 1 关注学习算法中的可学习组件：网络架构、初始参数、学习率等都可以被视为 $\phi$ 的一部分；不同元学习方法的差异，常常来自“把哪些组件放进 $\phi$”。}
{00:11:45--00:12:00}

常见可学习组件包括：
\begin{itemize}
    \item \textbf{初始参数：}学习一个好的起点，使新任务只需少量梯度更新即可适应。
    \item \textbf{学习率或更新步长：}不同层、不同参数甚至不同时间步可以有不同更新强度。
    \item \textbf{网络结构：}学习哪些模块、连接或结构适合一类任务。
    \item \textbf{距离或相似度度量：}在 few-shot classification 中，学习如何比较 support example 和 query example。
    \item \textbf{优化规则：}把“如何根据梯度更新参数”本身做成可学习程序。
    \item \textbf{任务表示：}学习如何把一个 task 的 support set 编码成可用于预测 query 的信息。
\end{itemize}

所以 $F_\phi$ 不是唯一形态。若 $\phi$ 是初始化参数，就得到 learning to initialize 的路线；若 $\phi$ 是度量函数，就接近 learning to compare；若 $\phi$ 是优化器，就接近 learning to optimize。课程此处先不进入具体算法，而是强调统一框架：这些方法都可以看成在定义一个带未知组件的学习算法。

\begin{warningbox}{容易混淆的两种参数}
    $\phi$ 不是普通分类器的参数。分类器参数通常由某个任务内部的训练资料产生；$\phi$ 则来自许多 training tasks 的经验，用来决定以后遇到新任务时该怎样快速产生分类器。简单说，$\theta$ 属于 task 内部，$\phi$ 属于跨 task 的学习经验。
\end{warningbox}

\subsection{本章小结}

元学习的 Step 1 是决定学习算法 $F_\phi$ 的候选空间。它可以学习初始化、学习率、结构、度量、优化规则或任务表示。不同元学习分支的本质差异，往往就是 $\phi$ 代表的东西不同。

% =====================================================================
\section{Step 2：如何评价一个学习算法}
% =====================================================================

定义了 $F_\phi$ 之后，第二步要定义损失函数 $L(\phi)$。在普通机器学习里，loss 衡量某个模型在训练样本上的预测错误；在元学习里，loss 衡量某个学习算法在一批任务上的表现。也就是说，评价对象从“一个模型”变成“一个会产生模型的算法”。

设有 $N$ 个 training tasks：
\[
    \Task_1,\Task_2,\ldots,\Task_N.
\]
每个任务都有自己的训练部分和测试部分：
\[
    \Task_n = \left(\Dtrain_n,\Dtest_n\right).
\]
对第 $n$ 个任务，学习算法先看该任务的训练资料：
\[
    f_n = \flearner(\Dtrain_n).
\]
然后用该任务的测试资料评价 $f_n$：
\[
    \ell^n(\phi)
    =
    \sum_{(x,y)\in \Dtest_n}
    \operatorname{CE}\bigl(f_n(x),y\bigr).
\]
最后把所有 training tasks 的评价加总：
\[
    L(\phi)
    =
    \sum_{n=1}^{N}\ell^n(\phi).
\]

符号说明如下。
\begin{description}
    \item[$\ell^n(\phi)$] 第 $n$ 个 task 上，学习算法 $F_\phi$ 产出的模型在该 task 测试部分上的错误。
    \item[$L(\phi)$] 跨所有 training tasks 的总损失，用来评价学习算法整体好坏。
    \item[$N$] 用来训练学习算法的任务数量，而不是单一任务里的样本数量。
\end{description}

\framefig{0.84}{figures/fig_05_meta_loss_across_tasks.jpg}
{Step 2 的核心：同一个学习算法 $F_\phi$ 会在多个 training tasks 上分别产出分类器，再用各 task 的测试部分计算 loss；所有 task loss 加总后才是 $L(\phi)$。}
{00:19:45--00:20:00}

这一步很反直觉：在每个 training task 内部，计算 meta-loss 时确实会用到这个 task 的 testing examples。原因是我们不是在训练该 task 的最终分类器，而是在训练学习算法；training task 的 testing examples 正好提供“这个学习算法是否真的会学”的评价信号。

\begin{importantbox}{为什么 training task 可以有 testing examples}
    对元学习而言，training task 整体都是训练资料。它内部再切成 train/test，是为了模拟未来新任务的使用方式：学习算法只能先看少量 train examples，然后必须在未见过的 examples 上表现好。因此，training task 的 test split 用来训练 $F_\phi$；真正不能碰的是最终 testing task 的资料。
\end{importantbox}

\begin{warningbox}{不要把 task 的 test split 和最终测试集混为一谈}
    “Task 内部的 testing examples”与“整个 meta-learning 实验的 testing tasks”是两层概念。前者在 meta-training 中可以用于计算 $L(\phi)$；后者用于评估最终学到的学习算法，训练时不能使用。
\end{warningbox}

\subsection{本章小结}

元学习的 Step 2 是定义学习算法的 loss。它不是看某个固定分类器在一个数据集上错多少，而是看同一个学习算法在许多 training tasks 中，能否用少量训练样本产生能泛化到该 task 测试部分的模型。

% =====================================================================
\section{Step 3：优化学习算法本身}
% =====================================================================

第三步是找出让 $L(\phi)$ 最小的学习算法参数：
\[
    \phi^\ast
    =
    \operatorname*{arg\,min}_{\phi} L(\phi).
\]
如果 $L(\phi)$ 对 $\phi$ 可微，就可以使用熟悉的梯度下降：
\[
    \phi \leftarrow \phi - \eta \nabla_{\phi}L(\phi).
\]
若 $L(\phi)$ 不可微，仍然可以把它当作黑箱优化问题，使用强化学习、演化算法或其他 derivative-free optimization 方法。

\framefig{0.84}{figures/fig_06_meta_optimization_step3.jpg}
{Step 3 把学习算法参数 $\phi$ 当作优化对象：若能计算 $\partial L(\phi)/\partial\phi$，就可用梯度下降；若不可微，也可考虑强化学习或演化算法等黑箱优化方法。}
{00:23:30--00:23:45}

优化完成后得到 $F_{\phi^\ast}$，也就是课程中的 learned learning algorithm。未来遇到一个新任务时，不再从手工学习算法开始，而是使用 $F_{\phi^\ast}$：
\[
    f_{\mathrm{new}}
    =
    F_{\phi^\ast}(\Dtrain_{\mathrm{new}}).
\]
如果元学习成功，$F_{\phi^\ast}$ 会把过去 training tasks 中积累的经验压缩进 $\phi^\ast$，使新任务适应更快、所需标注更少，或超参数搜索更少。

\begin{knowledgebox}{元学习训练本身可能很贵}
    元学习常常把“许多 task 内部的训练过程”嵌进一个更外层的优化中。即使每个 task 都很小，总成本也可能很高，因为每次更新 $\phi$ 都要模拟若干任务上的学习与评估。这也是为什么元学习论文常特别关注效率、近似梯度、任务采样和 inner loop 长度。
\end{knowledgebox}

\subsection{本章小结}

元学习的 Step 3 是在 task 层级上优化 $F_\phi$。普通机器学习优化 $\theta$，得到一个任务模型；元学习优化 $\phi$，得到一个可复用于新任务的学习算法。只要能定义 $L(\phi)$，就可以用梯度法或黑箱优化法寻找 $\phi^\ast$。

% =====================================================================
\section{Few-shot 框架中的任务、资料与循环}
% =====================================================================

课程接着把抽象框架落到 few-shot learning 语境。Few-shot learning 是目标：希望模型只看少量标注样本就能解决新任务。Meta learning 是常见手段：通过许多 training tasks 训练一个学习算法，使它在未来 testing task 上只需少量样本就能适应。

\framefig{0.84}{figures/fig_07_meta_framework_training_testing_tasks.jpg}
{元学习框架中的 training tasks 与 testing task：上方 training tasks 用来学习 $F_{\phi^\ast}$；真正关心的是左下方新的 testing task，它只提供少量 labeled training data，并要求模型在该 task 的 test examples 上表现好。}
{00:27:45--00:28:00}

这个图中最重要的是边界：training tasks 与 testing task 可以来自同一类任务分布，但 testing task 本身不能参与训练 $F_\phi$。如果把 testing task 的 train split 也拿来训练 $\phi$，就会破坏 few-shot/meta-learning 的评估设定。

\subsection{Support set 与 query set}

为了避免“训练资料”和“测试资料”这两个词在不同层级上混乱，meta-learning 文献常把 task 内部的训练部分叫 support set，把 task 内部的测试部分叫 query set。

\framefig{0.84}{figures/fig_08_support_query_sets.jpg}
{在元学习文献中，一个 task 内部常被切成 support set 与 query set：support set 用来让学习算法产生该 task 的模型，query set 用来计算这个 task 对学习算法的评价。}
{00:30:45--00:31:00}

这个命名非常有用：
\begin{itemize}
    \item \textbf{Support set：}给学习算法看的少量标注样本，负责“支持”当前 task 的适应。
    \item \textbf{Query set：}当前 task 中用于评价的样本，负责检验由 support set 产生的模型是否泛化。
    \item \textbf{Training tasks：}许多带 support/query 切分的任务，用来训练 $F_\phi$。
    \item \textbf{Testing task：}最终评估用的新任务，只能用其 support set 做适应，再在 query/test 上评估。
\end{itemize}

\begin{warningbox}{命名不统一会导致的误解}
    如果说“用 testing examples 训练”，听起来像数据泄漏；但如果说“用 training task 的 query set 计算 meta-loss”，含义就清楚得多。元学习中最危险的混乱，往往不是公式难，而是层级命名没有区分。
\end{warningbox}

\subsection{Within-task 与 across-task}

课程进一步区分两种 training：
\begin{itemize}
    \item \textbf{Within-task training：}在单个 task 内，根据 support set 得到该 task 的模型。
    \item \textbf{Across-task training：}跨许多 training tasks，学习出一个更好的学习算法 $F_\phi$。
\end{itemize}

\framefig{0.84}{figures/fig_09_within_across_task_training.jpg}
{普通机器学习只包含一个 task 内部的训练；元学习同时包含 across-task training 和 within-task training：先跨任务学习学习算法，再在每个新任务内部用少量 support examples 产生分类器。}
{00:32:00--00:32:15}

这一区分解释了元学习为什么经常看起来像“双层训练”。内层是 task 内部的模型适应，外层是学习算法本身的更新。对于 learning to initialize 类方法，内层可能是几步梯度下降，外层则根据 query loss 更新初始化参数。

\subsection{Loss 层级：样本求和与任务求和}

普通机器学习的 loss 往往是在一个任务的训练样本上求和：
\[
    L(\theta)
    =
    \sum_{k=1}^{K} e_k.
\]
元学习则多一层任务求和：
\[
    L(\phi)
    =
    \sum_{n=1}^{N}\ell^n,
    \qquad
    \ell^n
    =
    \sum_{m=1}^{M_n} e_m^n.
\]
其中 $e_m^n$ 是第 $n$ 个 task 的第 $m$ 个 query example 的错误，$\ell^n$ 是一个 task 的 query loss，$L(\phi)$ 是跨 training tasks 的总损失。

\framefig{0.84}{figures/fig_10_loss_machine_vs_meta.jpg}
{普通机器学习的 loss 对单一任务的训练样本求和；元学习的 loss 先在每个 task 的 query examples 上求和得到 $\ell^n$，再跨 training tasks 求和得到 $L(\phi)$。}
{00:34:45--00:35:00}

这个层级也说明为什么元学习需要 task distribution。若只有一个任务，就很难判断一个学习算法是否真的“会学习”；它可能只是记住了一个任务的特殊性。多个 tasks 提供了泛化压力，让 $F_\phi$ 学到跨任务可复用的规律。

\subsection{Episode 与 inner/outer loop}

在 few-shot 元学习中，一次“从 support set 适应，再在 query set 上评价”的流程常被称为一个 episode。Across-task training 会采样许多 episodes，用它们来更新 $F_\phi$。

\framefig{0.84}{figures/fig_11_inner_outer_loop.jpg}
{在需要显式模拟任务内训练的元学习方法中，常把 across-task training 称为 outer loop，把 within-task training 称为 inner loop；但这个叫法最适合 learning to initialize，不一定适用于所有元学习方法。}
{00:37:15--00:37:30}

课程特别提醒：outer loop 与 inner loop 的说法并非所有元学习方法都必须采用。它最自然地出现在 learning to initialize 一类方法中，因为这些方法确实有一个反复执行的任务内优化过程。但如果学习算法不是迭代式的，例如一次前向传播就输出分类器或度量结果，那么 inner loop 的概念就不一定恰当。

\begin{importantbox}{把 few-shot episode 看成一次模拟考试}
    Support set 相当于给学生看的复习资料，query set 相当于同一个 task 的考试题。Meta-training 不断采样许多这样的模拟考试，让学习算法学会“怎样利用少量资料准备考试”。最终 testing task 则是全新的考试，训练时不能提前看到。
\end{importantbox}

\subsection{本章小结}

Few-shot 语境下，元学习最容易混乱的是层级。Task 内部有 support/query；task 之间有 training/testing；训练过程又有 within-task 与 across-task。理解这些层级后，$L(\phi)$ 的结构就很自然：它衡量学习算法在许多任务中“看 support、答 query”的整体能力。

% =====================================================================
\section{总结与延伸}
% =====================================================================

课程最后强调：你在普通机器学习中熟悉的许多概念，通常也能迁移到元学习，只是层级要上移。普通机器学习担心 overfitting on training examples；元学习担心 overfitting on training tasks。普通机器学习可以增加训练样本；元学习可以增加 training tasks。普通机器学习有 data augmentation；元学习也可以做 task augmentation。普通机器学习需要开发集调超参数；元学习也可以设置 development tasks 来选择算法与超参数。

\framefig{0.84}{figures/fig_12_meta_vs_ml.jpg}
{Meta Learning 与普通 ML 的类比：普通 ML 中关于过拟合、增加训练数据、数据增强、超参数与开发集的经验，通常可以上移到 task 层级，在 meta-learning 中变成 training tasks、task augmentation 与 development tasks。}
{00:38:30--00:38:45}

把 P55 压缩成一句话：元学习不是要推翻机器学习，而是把机器学习的三步骤递归地用在“学习算法”上。普通机器学习的问题是：
\[
    \text{如何从资料中找出好函数？}
\]
元学习的问题是：
\[
    \text{如何从许多任务中找出一个好学习算法？}
\]

因此，本讲的核心对应关系可以整理如下。
\begin{center}
\begin{tabular}{L{0.27\textwidth} L{0.31\textwidth} L{0.31\textwidth}}
\toprule
概念 & 普通机器学习 & 元学习 \\
\midrule
待学习对象 & 任务函数 $f_\theta$ & 学习算法 $F_\phi$ \\
输入资料 & 一个任务的训练样本 & 多个 training tasks \\
输出结果 & 一个具体模型 $f_{\theta^\ast}$ & 一个 learned learning algorithm $F_{\phi^\ast}$ \\
评价方式 & 模型在样本上的 loss & 学习算法在多个 task query sets 上的 loss \\
泛化目标 & 泛化到同一任务的新样本 & 泛化到新 task \\
\bottomrule
\end{tabular}
\end{center}

\begin{importantbox}{P55 最重要的概念压缩}
    元学习的“资料”不是单个样本，而是一批任务；元学习的“模型”不是单个分类器，而是学习算法；元学习的“泛化”不是只到新样本，而是到新任务。把这三点抓住，后续看 MAML、metric-based meta-learning、learning to optimize 等方法时，就能知道它们究竟在学习哪一层。
\end{importantbox}

从后续学习角度看，本讲提供的是统一语法。之后遇到具体方法时，可以逐一追问四个问题：
\begin{enumerate}
    \item 它把学习算法 $F_\phi$ 里的哪一部分设为可学习？
    \item 它如何构造 task、support set 与 query set？
    \item 它的 meta-loss $L(\phi)$ 怎样计算，是否需要反传穿过 inner loop？
    \item 它最终希望在什么类型的新 task 上快速适应？
\end{enumerate}

这些问题比直接背方法名更有用。因为元学习方法很多，但它们大多都能放回 P55 的三步骤：定义 $F_\phi$，定义 $L(\phi)$，优化 $\phi$。

\subsection{本章小结}

P55 建立了元学习的基本框架：普通机器学习找函数，元学习找学习算法；普通 loss 在样本层级求和，meta-loss 在 task 层级求和；普通训练关注一个任务内的模型参数，meta-training 关注跨任务可复用的学习经验。后续所有具体元学习算法，都可以先放进这个框架中定位。

\end{document}
